\chapter{ADRs: registros de decisión}
\label{sec:adrs}

Este apéndice recoge las siete ADRs (Architectural Decision
Records) que fundamentan el diseño actual. El formato es el mismo para
todas: contexto breve, decisión, consecuencias. El capítulo
\ref{sec:decisiones} del cuerpo principal narra cómo encajan entre
sí; aquí quedan en su forma de referencia.

\textbf{Fuente primaria}: \filepath{docs/_archive/adrs/001..006.md}
en el repositorio. Si hay discrepancia entre este apéndice y esos
archivos, prevalece el archivo.

\section*{ADR-001 — Canonical API Layer}

\textbf{Fecha}: 2026-04-23. \textbf{Estado}: aprobada e
implementada (Fase 2).

\textbf{Contexto}. \filepath{platform/api} dependía de
\filepath{apps/cli} como librería para reusar
\code{sentinel_streaming}, creando un ciclo lógico
API $\to$ CLI $\to$ core. \filepath{apps/cli} además saltaba la API
yendo directo al core, bifurcando la ruta de ingesta.

\textbf{Decisión}. (D1) Extraer \code{sentinel_streaming} a un
crate nuevo \filepath{platform/siem} del que dependen API y CLI.
(D2) \filepath{apps/cli} pasa a ser cliente in-proc de
\filepath{platform/api} (descartada la variante HTTP). (D3) DTOs
en \filepath{platform/api/src/dto/v1/} versionados con semver.

\textbf{Consecuencias}. Cero ciclos en \code{cargo metadata}. Una
única lógica, varios transportes. Extender la API requiere versionar
DTOs y dejar el endpoint viejo funcional durante un ciclo.

\section*{ADR-002 — FFI Isolation}

\textbf{Fecha}: 2026-04-23. \textbf{Estado}: aprobada e
implementada (Fase 3).

\textbf{Contexto}. Un único \code{extern "C"} (bueno), pero RAII
parcial, invariante ``válido tras \code{finalize()}'' no codificada,
cinco funciones C sin uso desde Rust, cero tests de equivalencia,
build script sin chequeo de ABI.

\textbf{Decisión}. (D1) Extraer la fachada FFI a
\filepath{core/matcher-ffi}. (D2) Typestate para \code{GmGraph}
con \code{Building}/\code{Finalized}. (D3) RAII completo en
\code{GmMatcher}. (D4) \code{gm_abi_version()} con check en
load-time; mismatch $\Rightarrow$ crash temprano. (D5) Property
tests de equivalencia DFS $\leftrightarrow$ SIMD.

\textbf{Consecuencias}. El aislamiento FFI pasa de convención a
garantía de tipo. Imposible ejecutar un hunt sobre un grafo no
finalizado. Imposible shippear con un \filepath{libgraphmatch.so}
de ABI distinta sin que el binario crashee al arranque.

\section*{ADR-003 — GNN v2 Experimental}

\textbf{Fecha}: 2026-04-23. \textbf{Estado}: aprobada; corpus
export y gate implementados; modelo no entrenado para Fortigate.

\textbf{Contexto}. \filepath{gnn_v2.onnx} referenciado en
producción por string literal, no entrenado sobre todos los
datasets objetivo. Binario enviado con skeleton produce
probabilidades sin fallar visiblemente.

\textbf{Decisión}. (D1) Tres componentes separados:
\filepath{core/gnn/v1} estable, \filepath{core/gnn/v2} detrás de
\code{feature = "gnn-v2-experimental"}, \filepath{core/gnn/common}
con el gate. (D2) \code{ModelGate::validate} rechaza archivos no
firmados o con dimensiones incorrectas. (D3) Corpus export
(\filepath{platform/sources/aad_training_export.rs}) separado del
inferencer.

\textbf{Consecuencias}. Enviar un binario con GNN v2 no validado
requiere (a) flippear el feature flag y (b) pasar el gate.
Accidentalmente imposible. Extender a Fortigate exige nuevo corpus
export + entrenamiento.

\section*{ADR-004 — MCP Tools Architecture}

\textbf{Fecha}: 2026-04-23. \textbf{Estado}: aprobada e
implementada (Fase 4).

\textbf{Contexto}. \filepath{graph-hunter-mcp/src/index.ts} era un
monolito de \raisebox{0pt}[0pt][0pt]{\textasciitilde}1670 líneas con 54 tools hardcoded, sin
contrato tipado entre definición de tool y endpoint HTTP.

\textbf{Decisión}. (D1) Interfaz \code{Tool<I, O>} tipada con
\code{stability} y \code{version}. (D2) Registry modular por
dominio (10 dominios). (D3) Un archivo por tool. (D4) Tests
contractuales contra \filepath{platform/api}. (D5) El paquete vive
en \filepath{platform/mcp/}; el paquete legacy
\filepath{graph-hunter-mcp/} queda un ciclo más por compat.

\textbf{Consecuencias}. Cambiar una tool es un PR localizado.
Deprecar una tool es una operación gobernable. Mismatches de
contrato se ven en tests, no en producción.

\section*{ADR-005 — DSL Extensibility}

\textbf{Fecha}: 2026-04-23. \textbf{Estado}: aprobada e
implementada (Fase 3).

\textbf{Contexto}. El DSL vivía en
\filepath{core/graph_hunter_core/src/dsl.rs}, con entity types y
relation types hardcoded. Clientes pedían extensiones
(\code{Container}, \code{ServiceAccount}, predicados temporales).

\textbf{Decisión}. (D1) Extraer a \filepath{platform/dsl}. (D2)
Trait \code{DslExtension} para registrar entity/relation/predicate.
(D3) Versionar gramática: \code{v1} estable, \code{v2-experimental}
detrás de feature. (D4) Preservar compatibilidad de hipótesis
escritas contra v1.

\textbf{Consecuencias}. El DSL deja de ser código del núcleo y pasa
a ser superficie extensible. Un cliente puede extender sin bloquear
a otros. Versiones futuras viven en paralelo hasta promoción.

\section*{ADR-006 — WIP Absorption}

\textbf{Fecha}: 2026-04-23. \textbf{Estado}: aprobada e
implementada (Fase 2, commit 0).

\textbf{Contexto}. \code{git status} al arrancar Fase 0: 44
archivos modificados + 14 untracked. No era código abandonado:
iteraciones M4 (agentic slow-lane), M5 (DLQ), M6 (invariant
checker) coherentes con el rediseño. Ignorarlas habría perdido
trabajo; separarlas antes de reorganizar habría diluido la señal de
cada commit.

\textbf{Decisión}. Tres commits atómicos \emph{antes} de la
reorganización: (1) M4 agentic + local-llm, (2) M5 DLQ, (3) M6
invariants. Cada commit incluye código + tests + DTOs asociados.

\textbf{Consecuencias}. El rediseño arranca desde una base
coherente. Los commits siguientes son puramente estructurales, no
mezcla de nuevas features con cambios arquitectónicos.

\section*{ADR-007 — Ingestor schema-agnóstico en tres carriles}
\label{sec:adr-007}

\textbf{Fecha}: 2026-04-24. \textbf{Estado}: aprobada, en
implementación (fast lane en \filepath{platform/vrl}; slow/learning
lanes en \filepath{platform/api/src/operations/agentic*.rs}).

\textbf{Contexto}. La ingesta heurística acopla un parser por
formato en el binario del daemon. Cada formato nuevo o cambio
aguas arriba (campo agregado en Sentinel, separador nuevo en
FortiOS) obliga a release para que el daemon lo procese. El parser
generic cubre casos conocidos y falla silencioso en los demás.
Resultado: un SOC con un formato no reconocido se queda sin ingesta
utilizable, o pierde fidelidad con el parser generic. No hay
camino incremental: ninguno de los parsers aprende de su propia
operación.

\textbf{Decisión}.

\begin{enumerate}
    \item \textbf{Fast lane}: separar la ingesta caliente en un
    programa VRL compilado determinísticamente desde un
    \code{FieldConfig}. Cero LLM en el camino crítico.
    \item \textbf{Slow lane}: usar un LLM \emph{offline} como
    compilador de muestras a \code{FieldConfig}. El output del LLM
    es datos, no código. Todo output pasa por revisión humana antes
    de entrar a producción.
    \item \textbf{Learning lane}: el fast lane emite métricas
    estructurales por batch; un monitor de drift dispara re-review
    en el slow lane sin parar la ingesta.
    \item \textbf{Mapping library}:
    \filepath{platform/mapping-library} persiste los mappings
    aprobados con \code{mapping_hash = sha256(VRL)} como identidad
    estable y provenance-stamp por evento ingestado.
    \item \textbf{Target canonical}: los mappings emiten OCSF v1.4
    como shape canonical. El parser no decide qué es una entidad;
    el FieldConfig-a-OCSF sí.
\end{enumerate}

\textbf{Consecuencias}.

\begin{itemize}
    \item La frontera humano/máquina queda bien puesta: LLM
    propone, humano aprueba, compilador emite, runtime ejecuta. Cada
    paso es trazable.
    \item El daemon no necesita release por formato nuevo: se
    publica un nuevo \code{FieldConfig} y el runtime lo recoge.
    \item La determinística del fast lane se preserva:
    \code{mapping_hash} estable, reproducible byte-a-byte.
    \item Superficie nueva: tools MCP
    (\code{parser_generate}, \code{publish_mapping},
    \code{agentic_review_*}, \code{schema_drift_detect},
    \code{mapping_regression_test}, \code{fetch_shared_mappings})
    que juntas materializan el carril agéntico.
    \item Dependencia operativa nueva (backend LLM). Mitigada por
    intercambiabilidad vía \filepath{platform/local-llm} y por que
    el carril es offline — caída del LLM no afecta ingesta viva.
\end{itemize}

\textbf{Alternativas consideradas}.

\begin{itemize}
    \item \textbf{LLM online en el camino crítico}: descartado por
    latencia, costo y superficie de falla. El LLM puede parsear mal
    un batch productivo; el humano, no. Y el costo por evento
    ingerido sería prohibitivo.
    \item \textbf{Parser genérico mejorado}: descartado porque
    hereda la misma falla — no aprende de su propia ejecución ni se
    adapta a cambios aguas arriba.
    \item \textbf{Schema-registry estático aportado por usuario}:
    factible pero exige que el usuario escriba FieldConfig a mano.
    El slow lane automatiza ese paso sin perder el gate humano.
\end{itemize}
